%%%%%%%%%%%%%%%%%%%%%%%%%%%%%%%%%
% Introduction                  %
%%%%%%%%%%%%%%%%%%%%%%%%%%%%%%%%%

\section{INTRODUCTION}
\label{sec:introduction}

This section provides an overview of the conducted research, starting with the project context, followed by the purpose of the study, research objectives, research questions, and the approach and boundaries that define the scope of this work.

\subsection{Project Context}

Sound synthesis is a fundamental component of modern music production and cinematic sound design, enabling the creation of complex, rich sonic textures as well as faithful emulations of acoustic instruments and sounds found in nature. While synthesis can be achieved through purely acoustic means, synthesizers have become one of the primary tools for this purpose. At their core, synthesizers expose a structured interface of parameters, including oscillators, filters, envelopes, modulation routings, and effects, whose combined interaction shapes the output sound. The relationship between these parameters and the resulting audio is highly nonlinear and unintuitive, making manual sound reproduction a time-consuming process \cite{roads1996computer, cook2002realsound}.

\par \vspace{0.5cm}

\noindent Historically, the most culturally influential synthesizers, namely the Minimoog by Moog and the Prophet-5 by Sequential Circuits, featured a fixed signal flow with a comparatively small number of parameters. Despite this simplicity, these instruments remain widely used in modern music production and are capable of expressing a surprisingly wide range of sounds. This stands in contrast to modern software synthesizers such as Vital by Matt Tytel \cite{tytel2020vital}, Serum by Xfer Records, or Diva by u-he, which expose hundreds of parameters, many of which remain unexplored in practice. Notably, even when working with these feature-rich instruments, producers tend to gravitate toward a familiar and interpretable subset: oscillator waveforms, envelope shapes, filter cutoff, and a few modulation amounts. This observation led us to restrict the parameter space in this study to a curated subset inspired by those vintage signal chains, keeping the task both musically meaningful and tractable.

\par \vspace{0.5cm}

\noindent In this project, we focus on Vital, a modern open-source wavetable synthesizer developed by Matt Tytel and widely adopted in the music production community, making it an ideal subject for this research. Our motivation is straightforward: a producer who hears a sound they want to recreate, whether from a reference track or their own imagination, typically faces a time-consuming process of manual programming. An automated estimation system would serve as a shortcut for beginners and a time-saving tool for experienced sound designers alike. More broadly, it can act as a creative aid: feeding an arbitrary audio signal into the estimator yields a synthesizer approximation that can open timbral directions difficult to discover through manual exploration.

\par \vspace{0.5cm}

\noindent The task of recovering synthesizer parameters from audio, known as inverse synthesis, has been studied since the early days of computational audio. Initial approaches relied on Genetic Algorithms (GA) \cite{walker2007audiosynthesis}, which are computationally expensive and scale poorly to high-dimensional parameter spaces. Deep learning later shifted the paradigm toward direct regression \cite{yeeking2018automatic, barkan2019inversynth}, removing the need for iterative synthesis at inference time. More recently, architectures from computer vision have been adapted for audio: Convolutional Neural Networks (CNNs) on mel-spectrograms have become a standard baseline \cite{hershey2017cnn}, and the Audio Spectrogram Transformer (AST) has shown strong results using global self-attention \cite{gong2021ast}. The work most directly related to this project is that of Bruford et al. \cite{bruford2024ast}, who applied AST to synthesizer sound matching and established a strong transformer baseline for the task.

\subsection{Purpose of the Study}

The purpose of this work is to evaluate how effectively existing deep learning architectures can estimate synthesizer parameters directly from audio, comparing them under a shared dataset and experimental setup. To support this, we design and implement a modular training and evaluation framework applicable to multiple architectures and training configurations, using the Vital synthesizer throughout. Audio samples are generated programmatically by driving Vital through the Pedalboard library under controlled rendering conditions; each sample is paired with its ground-truth parameter vector and the Musical Instrument Digital Interface (MIDI) pitch at which it was rendered. The dataset is built in two stages: an initial generated version and a subsequently augmented version three times larger, obtained by rendering each preset at multiple MIDI pitches. Across three experimental iterations of increasing complexity, we compare three neural architectures: a residual CNN baseline \cite{he2016resnet} (first iteration only), a ConvNeXt backbone \cite{liu2022convnext}, and an AST \cite{gong2021ast}, tracing the contributions of multi-scale mel input, Feature-wise Linear Modulation (FiLM)-based pitch conditioning \cite{perez2018film}, and multi-pitch data augmentation to overall estimation accuracy.

\subsection{Research Objectives}

The following objectives define the scope of this work:

\begin{itemize}

\item[$\blacksquare$] \textbf{Objective 1:} Evaluate and compare three deep learning architectures, namely a residual CNN baseline, ConvNeXt, and AST, in their ability to regress synthesizer parameters from mel-spectrogram inputs, and identify which architectural family is best suited to this task.

\item[$\blacksquare$] \textbf{Objective 2:} Assess whether a multi-scale mel-spectrogram input combining fine, mid, and coarse time-frequency resolutions provides a measurable improvement over a single-scale mel-spectrogram for synthesizer parameter estimation.

\item[$\blacksquare$] \textbf{Objective 3:} Investigate the conditions under which FiLM-based pitch conditioning improves parameter estimation accuracy, and determine how the availability of multi-pitch training data affects its effectiveness.

\end{itemize}

\subsection{Research Questions}

\begin{itemize}

    \item[$\blacksquare$] \textbf{Research Question 1:} How do convolutional architectures (residual CNN and ConvNeXt) compare to AST in regressing synthesizer parameters from mel-spectrograms? \\
    \textbf{Rationale:} CNNs build up representations by stacking local convolutions, while transformers reason globally through self-attention across the entire input. These are fundamentally different inductive biases, and it is not obvious which one better serves a synthesizer parameter regression task. Synthesized sounds tend to be characterized by localized spectral structures such as harmonic stacks, filter shapes, and envelope transients, which intuitively favor local feature extraction. Whether this holds in practice, and by how much, is what this comparison aims to settle.

    \item[$\blacksquare$] \textbf{Research Question 2:} Does replacing a single mel-spectrogram with a three-scale representation combining fine, mid, and coarse spectrograms improve synthesizer parameter estimation accuracy? \\
    \textbf{Rationale:} Any single mel-spectrogram is a resolution compromise: you cannot simultaneously maximize temporal and spectral precision. This matters because synthesizer parameters operate at different timescales: a fast attack is best seen in a fine-resolution spectrogram, while harmonic texture and filter shape are more legible at coarser frequency resolution. Providing all three scales as a stacked input is a practical way to sidestep this trade-off, and this question tests whether doing so actually translates to better parameter estimates.

    \item[$\blacksquare$] \textbf{Research Question 3:} Under what conditions does FiLM-based pitch conditioning improve synthesizer parameter estimation, and why does its effectiveness depend on the training data? \\
    \textbf{Rationale:} The frequency content of a synthesized note depends on both its synthesis parameters and the MIDI pitch at which it is played. A model trained on single-pitch renders never sees the same preset at different pitches, so there is no signal pushing it to make use of the pitch input. The hypothesis is that multi-pitch rendering, where each preset is rendered at several MIDI notes, forces the model to learn pitch-invariant parameter representations and gives the FiLM conditioning a meaningful role. This question tests that hypothesis directly.

\end{itemize}

\subsection{Approach and Boundaries of the Study}

This research is structured around the Design Science Research (DSR) methodology \cite{wieringa2014design}. DSR is a problem-solving paradigm centered on the creation and evaluation of purposeful artifacts; here, those artifacts are the deep learning models and the data generation pipeline built to support them. Its iterative cycle of problem investigation, design, implementation, and evaluation fits naturally with how this project developed: each experimental iteration (V1, V2, V3) refined the system based on what the previous one revealed.

\subsubsection{Scope}
Three neural architectures are compared: a residual CNN serving as a V1 reference, ConvNeXt, and AST, all applied to synthesizer parameter regression on the Vital synthesizer. The comparison covers single-scale and multi-scale mel-spectrogram inputs, FiLM-based pitch conditioning, grouped prediction heads, and multi-pitch data augmentation across three experimental iterations. We report Mean Absolute Error (MAE) in the normalized parameter space, with 5-fold cross-validation and a per-parameter breakdown for each configuration. For V2 and V3, a Multi-Scale Spectral (MSS) loss computed on re-rendered audio serves as an additional perceptual evaluation metric.

\subsubsection{Limitations}
This investigation ran into several practical walls. GPU access was limited to a single NVIDIA RTX 4090 shared across multiple researchers; scheduling conflicts meant we could not run as many training experiments as we would have liked. The data pipeline added another layer of constraint: real-time audio rendering through Vital generated roughly 140 GB of raw Waveform Audio File Format (WAV) data, and the sheer time cost of rendering and validating every clip put a hard ceiling on how large the dataset could grow. The inverse synthesis literature also proved to be a narrow field, with few papers close enough to this work to serve as direct benchmarks. On top of that, the project sits at an unusual intersection of audio signal processing, music theory, and deep learning, and getting up to speed across all three required real effort from a team that came in primarily as software engineers.

\subsubsection{Delimitations}
The study is scoped to the Vital software synthesizer and the 16 parameters defined earlier. Generative architectures, recurrent networks, and reinforcement learning are out of scope; the focus is squarely on direct regression. Each preset is rendered as a 4-second audio clip. This duration was chosen because the target sounds are one-shots and sustained patches of the kind used in modern music production, where individual notes rarely extend beyond a few seconds. A 4-second window is sufficient to capture the full amplitude envelope of most patches, including pads with slow attacks and long decays. As a consequence of this fixed-duration render with no note-off event, the amplitude envelope sustain and release stages are never triggered and cannot be inferred from the audio; they are therefore excluded from the parameter set. Clips are stored in WAV rather than a lossy-compressed format such as Moving Picture Experts Group Audio Layer III (MP3), as compression introduces spectral artifacts that could corrupt the fine timbral detail the models are trained to predict. Although MSS loss has been used as a training objective in differentiable synthesis frameworks \cite{engel2020ddsp, caspe2022ddx7}, it is not used for training here. Vital is a non-differentiable black-box Virtual Studio Technology (VST) plugin, so computing MSS at each training step would require a differentiable proxy synthesizer, which is well beyond the scope of this work. MSS is therefore reserved as a post-hoc perceptual evaluation metric for V2 and V3. The study does not extend to analog hardware synthesizers, other VST plugins, or broader parameter spaces.

\subsubsection{Assumptions}
Three working assumptions underlie this study. First, the timbre-constrained generation strategy produces a dataset with sufficient coverage of musically meaningful sounds for the models to learn generalizable synthesis mappings. Second, the mel-spectrogram captures enough spectral and temporal information to characterize all 16 target parameters, even though phase information is discarded. Third, a reduction in normalized MAE is taken to correlate with a perceptually meaningful improvement in how closely the re-rendered output matches the target sound.

\subsection{Contributions}

This work makes the following contributions to the field of inverse synthesis and audio deep learning:

\begin{itemize}

\item[$\blacksquare$] \textbf{First direct comparison of ConvNeXt and AST for synthesizer parameter estimation.} Prior work has applied AST to sound matching \cite{bruford2024ast} or convolutional models to parameter regression \cite{barkan2019inversynth}, but no study has directly compared these two architectural families in the same controlled experimental setting.

\item[$\blacksquare$] \textbf{First use of explicit pitch conditioning for synthesizer parameter regression.} A MIDI pitch input is incorporated as an auxiliary conditioning signal via FiLM modulation, enabling the model to disentangle fundamental frequency from timbral characteristics. This conditioning strategy has not previously been explored for this task.

\item[$\blacksquare$] \textbf{Multi-scale mel spectrogram as a three-channel input representation.} To the best of our knowledge, stacking fine, mid, and coarse mel spectrograms into a single multi-channel input tensor has not previously been applied to synthesizer parameter estimation, and its benefit over single-scale input is evaluated through a dedicated ablation.

\item[$\blacksquare$] \textbf{Timbre-constrained dataset generation methodology.} A structured dataset generation pipeline is proposed that constrains parameter sampling to musically coherent timbral families, ensuring coverage of practically relevant sounds while avoiding silent or pathological outputs.

\item[$\blacksquare$] \textbf{Open-source dataset generation pipeline for VST synthesizers.} The full Pedalboard-based rendering pipeline, including timbre-constrained preset generation, silence filtering, and multi-scale mel extraction, is released as open-source code. The pipeline is designed to be extensible: new parameters can be added to the generation schema with minimal configuration, making it straightforward to expand the parameter space or adapt the pipeline to other VST synthesizers. To our knowledge, no prior inverse synthesis work has publicly released a comparable VST-driven dataset generation framework.

\item[$\blacksquare$] \textbf{Open-source training and evaluation framework.} The complete experimental codebase, including all model architectures, training configurations, ablation variants, and cross-validation splits, is made publicly available, enabling direct reproduction and extension of all results reported in this study.

\end{itemize}

\subsection{Structure}

The remainder of this report is organized as follows. Section 2 establishes the theoretical background and reviews prior work in sound synthesis, inverse synthesis methods, and the deep learning architectures used in this study. Section 3 describes the research methodology, detailing the DSR framework applied, the data generation pipeline, and the experimental setup. Section 4 presents the three experimental iterations — V1, V2, and V3 — describing the design decisions, architectural choices, and incremental improvements introduced at each stage. Section 5 reports the quantitative results and findings across all experiments. Section 6 discusses the results in light of the research questions, addresses limitations, and interprets per-parameter performance patterns. Section 7 concludes the report by summarizing contributions and outlining directions for future work.
